\section{A free six-DOF model and its numerical realization}
\label{gnc:sec:free-simulation}
Take the body origin at the center of mass, SI units, constant mass $m>0$
and constant physical inertia $J=J_B=J^T>0$. Its principal moments satisfy
$J_i\le J_j+J_k$. Write $\mathcal R:S^3\to SO(3)$ for our Hamilton,
scalar-first quaternion map, $q=(q_0,\mathbf q)$, with
$\mathcal R(q)\mathbf a_B=\mathbf a_I$. The state and evolution are
\begin{align}
 \mathcal X&=\mathbb R^3\times\mathbb R^3\times S^3\times\mathbb R^3,
 &z&=(\mathbf r_I,\mathbf v_I,q,\boldsymbol\omega_B),\notag\\
 \dot{\mathbf r}_I&=\mathbf v_I,
 &\dot{\mathbf v}_I&=m^{-1}\bigl(\mathcal R(q)\mathbf F_B+
                              \mathbf F_{I,\mathrm{ext}}\bigr),\notag\\
 \dot q&=\tfrac12 q\otimes(0,\boldsymbol\omega_B),
 &\dot{\boldsymbol\omega}_B
 &=J^{-1}\bigl(\boldsymbol\tau_B-
                  \boldsymbol\omega_B\times J\boldsymbol\omega_B\bigr).
 \label{gnc:simulation-rhs}
\end{align}
Thus $\dim\mathcal X=12$ although $z$ stores $13$ real components;
$q\sim-q$ gives the same physical state. The moment $\boldsymbol\tau_B$
is about the center of mass; a force at body offset $\boldsymbol\rho_B$
contributes $\boldsymbol\rho_B\times\mathbf F_B$.

For $\mathbf F_B=\mathbf F_{I,\mathrm{ext}}=\boldsymbol\tau_B=0$,
\begin{align}
 \mathbf r_I(t)&=\mathbf r_I(0)+t\mathbf v_I(0),
 &\mathbf p_I&=m\mathbf v_I=\mathrm{const},\notag\\
 \mathbf H_I&=\mathcal R(q)J\boldsymbol\omega_B=\mathrm{const},
 &E&=\tfrac12m\lVert\mathbf v_I\rVert^2+
            \tfrac12\boldsymbol\omega_B^TJ\boldsymbol\omega_B
            =\mathrm{const}.\label{gnc:simulation-invariants}
\end{align}
Here $\mathbf H_I$ is spin about the center of mass; total angular momentum
about the inertial origin is $\mathbf r_I\times\mathbf p_I+\mathbf H_I$.
Generally $\dot{\boldsymbol\omega}_B\ne0$ even with no external wrench.
For distinct principal moments, steady spin about the largest or smallest
principal moment is stable; spin about the intermediate moment is unstable
(Goldstein \S5.6, pp.~200--208 \cite{dgdt:GPS2001}).

\subsection{Midpoint momentum and a Cayley attitude step}
For step size $h>0$, let $\boldsymbol\ell_n=J\boldsymbol\omega_{B,n}$.
For a body-fixed wrench held constant over the step, solve
\begin{align}
 \bar{\boldsymbol\ell}&=\tfrac12(\boldsymbol\ell_n+
                                              \boldsymbol\ell_{n+1}),
 &\bar{\boldsymbol\omega}&=J^{-1}\bar{\boldsymbol\ell},\notag\\
 \boldsymbol\ell_{n+1}-\boldsymbol\ell_n
 &=h\bigl(\boldsymbol\tau_B-
               \bar{\boldsymbol\omega}\times\bar{\boldsymbol\ell}\bigr),
 &u&=\frac{(1,\tfrac h2\bar{\boldsymbol\omega})}
             {\sqrt{1+\tfrac{h^2}4\lVert\bar{\boldsymbol\omega}\rVert^2}},
 \notag\\
 q_{n+1}&=q_n\otimes u,
 &R_{n+1}&=R_n\,
 \underbrace{(I-\tfrac h2\widehat{\bar{\boldsymbol\omega}})^{-1}
             (I+\tfrac h2\widehat{\bar{\boldsymbol\omega}})}_
             {\operatorname{cay}(h\widehat{\bar{\boldsymbol\omega}})}.
 \label{gnc:midpoint-cayley}
\end{align}
The nonlinear solve uses its residual, rather than a fixed number of
iterations, to decide convergence. Let
$q_{1/2}=(q_n+q_{n+1})/\lVert q_n+q_{n+1}\rVert$ using these continuous
quaternion lifts. Then
\begin{align*}
 \mathbf a_{1/2}
 &=m^{-1}\bigl(\mathcal R(q_{1/2})\mathbf F_B+
                                      \mathbf F_{I,\mathrm{ext}}\bigr),\\
 \mathbf v_{n+1}&=\mathbf v_n+h\mathbf a_{1/2},\qquad
 \mathbf r_{n+1}=\mathbf r_n+\tfrac h2(\mathbf v_n+\mathbf v_{n+1}).
\end{align*}
In exact arithmetic, for zero wrench,
\begin{align*}
 T_{\mathrm{rot},n+1}-T_{\mathrm{rot},n}
 &=\bar{\boldsymbol\omega}\cdot
                 (\boldsymbol\ell_{n+1}-\boldsymbol\ell_n)=0,\\
 (I+\tfrac h2\widehat{\bar{\boldsymbol\omega}})
       \boldsymbol\ell_{n+1}
 &=(I-\tfrac h2\widehat{\bar{\boldsymbol\omega}})
       \boldsymbol\ell_n
 \ \Longrightarrow\ R_{n+1}\boldsymbol\ell_{n+1}
                         =R_n\boldsymbol\ell_n.
\end{align*}
Floating-point conservation is limited by the nonlinear residual and
roundoff. This is a second-order trajectory approximation, not the exact
flow: even spherical inertia gives the step angle
$2\arctan(h\lVert\boldsymbol\omega\rVert/2)$ rather than
$h\lVert\boldsymbol\omega\rVert$.
Check both invariant drift and attitude/position convergence against an
independent reference; invariant preservation alone cannot bound phase error.

\subsection{What the model resolves; what to read next}
Equations~\eqref{gnc:simulation-rhs} suffice for an ideal rigid body in
force-free, torque-free space. A browser clock requests elapsed time;
the backend advances fixed physical steps. Applied controls are explicit
external wrenches; releasing them returns to free coast.

For dynamics, read Goldstein \S\S5.1, 5.3--5.6 (pp.~184--208)
\cite{dgdt:GPS2001}; for an engineering connection, Wie
\S5.5 (p.~339) and \S\S6.1--6.7 (pp.~349--371)
\cite{BWie2008}. Wie \S\S7.3--7.7 add quaternion feedback, momentum
management and jet selection; Chapters~8--9 add flexibility and slosh.
Spacecraft fidelity additionally requires measured mass properties and
appropriate environment/actuator models. Numerical verification does not
by itself validate those physical assumptions. NASA's six-DOF check cases
\cite{NASA6DOF2015} provide a later independent benchmark target.
